\section{Activation-space analysis}
\label{sec:qwen-activation-projection}

We next analyze whether the behavioral conformity effect is reflected in Qwen's
residual stream. The prompts and social manipulations follow the base
multi-actor setup of \citet{zhu2025conformity}; the activation analysis examines
the resulting assistant-start states under these controlled social contexts.

\paragraph{Diffmean generation.}
For each selected layer $\ell$, we construct a direction that contrasts two
matched social contexts for the same item. The first context repeats the same
incorrect answer across the previous participants; the second keeps the
multi-participant format but replaces the repeated consensus with random
alternative answers. All activations are taken at the assistant-start position,
the token position immediately before the model begins generating its answer.
Let
$h^\ell_{i,\mathrm{wrong}}$ and $h^\ell_{i,\mathrm{random}}$ denote the
residual-stream activation at the assistant-start position for item $i$ under
these two contexts. The diffmean direction is
\[
  d^\ell =
  \frac{1}{N}\sum_{i=1}^{N}
  \left(h^\ell_{i,\mathrm{wrong}} - h^\ell_{i,\mathrm{random}}\right),
\]
with $N=100$ items for each tensor. The random-answer contrast preserves
much of the prompt length and participant-list structure, so the direction is
aimed at repeated agreement on the same wrong answer rather than the generic
presence of other participants. The main direction in this section is the MMLU
$n=8$ \texttt{all\_wrong\_minus\_random\_answers} residual direction. This is a
high-pressure MMLU condition with substantial behavioral conformity
(0.281 under base prompting), but it is not only the endpoint condition. We also
computed the analogous MMLU $n=10$ direction and use it as a check in the
cosine analysis below.

We use two activation summaries. A residual direction uses the accumulated
residual-stream vector at the assistant-start token after a layer. A delta
direction uses the update contributed by that layer at the same position,
i.e. the layer output that is added into the residual stream. The residual view
asks whether the social signal is present in the model state; the delta view is
a stricter check for whether a particular block's update carries a similar
signal.

\paragraph{Activation oracle responses.}
For the MMLU $n=8$ direction, we also asked an activation oracle for a short
description of each layer's diffmean vector. The oracle is the Qwen3-8B
LatentQA LoRA used by \texttt{layer\_sweep.py}. Each layer is queried with a
free-form text prompt, a free-form belief/feature prompt, and a binary prompt:
``Does this activation represent pressure toward an incorrect consensus?''

The residual MMLU $n=8$ direction gives a fairly coherent pattern. The binary
oracle prompt marks layers 10, 12, and 22--33 as positive, with the dense part
of the run in layers 22--33. The free-form responses in this band are also not
just generic answer-format descriptions. Layer 23 is described as group behavior
influencing the assistant's responses; layer 24 as actions influenced by the
majority; layers 26, 29, 30, and 33 as treating the majority as correct; and
layers 31--32 as conformity being important for group success. The residual
direction alone shows where the signal is linearly accessible in the model
state; the delta direction gives a more local view of where that signal is being
written. On this criterion, layers 22 and 23 are the key layers: both are
positive in the delta summary, and both sit inside the larger residual band.
For the cosine analysis, we therefore use layer 23. This choice treats layer
23 as a point where the conformity-related information is being made linearly
available in the residual stream, while avoiding a selection rule based only on
large late-layer norms.

\begin{table}[t]
  \centering
  \caption{Activation-oracle responses for the Qwen MMLU $n=8$ diffmean.
  ``True layers'' are layers where the binary oracle prompt says the direction
  represents pressure toward an incorrect consensus.}
  \label{tab:qwen-oracle-layers}
  \begin{tabular}{ll}
    \toprule
    Activation summary & True layers \\
    \midrule
    Residual stream &
      10, 12, 22--33 \\
    Layer update (delta) &
      8, 22, 23 \\
    \bottomrule
  \end{tabular}
\end{table}

The delta results are much sparser than the residual results: only layers 8, 22,
and 23 are marked positive. We interpret this as evidence that the residual
stream contains a broader conformity-related direction over layers 22--33, while
the layer updates at 22 and 23 are where this information is most clearly being
computed or written into a linearly readable form. This interpretation is still
descriptive: the oracle responses motivate a layer choice for cosine
comparison, but do not by themselves establish that the direction causally
mediates the behavioral effect.

\paragraph{Cosine similarity across diffmean directions.}
We next ask whether the layer-23 MMLU direction is stable across related
conditions. A condition $c$ is a fixed tuple of dataset, social condition
($n=2,\ldots,10$, QD, or DA), and prompt style (base or explanation prompting),
for example MMLU $n=6$ base or BBH $n=10$ explain. For two conditions $a$ and
$b$, with layer-23 diffmean directions $d^\ell_a$ and $d^\ell_b$, we compute
\[
  \mathrm{cos}(a,b;\ell) =
  \frac{d^\ell_a \cdot d^\ell_b}{\|d^\ell_a\|\,\|d^\ell_b\|},
  \qquad \ell=23.
\]
This is a comparison between contrast directions, not a projection of mean
condition activations. It therefore tests whether different social conditions
yield a similar repeated-wrong-minus-random-answer direction in the residual
stream.

Table~\ref{tab:qwen-diffmean-cosines} reports layer-23 cosine similarities
against the MMLU $n=8$ and $n=10$ base references. The MMLU directions are
stable across high-participant MMLU conditions: the MMLU $n=8$ reference has
cosines 0.961, 0.972, 0.953, and 0.932 with MMLU base $n=6,7,9,10$,
respectively, and the MMLU $n=10$ reference has cosine 0.978 with MMLU base
$n=9$. Explanation prompting preserves the broad direction but lowers the mean
cosine. BBH directions are also positively aligned after the low-participant
conditions, but the alignment is weaker than within MMLU. For example, BBH
$n=10$ base has cosine 0.388 with the MMLU $n=8$ reference and 0.542 with the
MMLU $n=10$ reference.

\begin{table}[t]
  \centering
  \caption{Layer-23 cosine similarity between MMLU reference diffmeans and
  numeric participant-count diffmeans. Means and correlations are over
  $n=2,\ldots,10$; the MMLU base row for the reference condition includes the
  self-cosine of 1.}
  \label{tab:qwen-diffmean-cosines}
  \begin{tabular}{lllrrr}
    \toprule
    Reference & Dataset & Prompt & Mean cos. & $r(n,\mathrm{cos})$ & Cos. at $n=10$ \\
    \midrule
    MMLU $n=8$ & MMLU & base & 0.769 & 0.730 & 0.932 \\
    MMLU $n=8$ & MMLU & explain & 0.695 & 0.793 & 0.899 \\
    MMLU $n=8$ & BBH & base & 0.430 & 0.436 & 0.388 \\
    MMLU $n=8$ & BBH & explain & 0.335 & 0.582 & 0.397 \\
    MMLU $n=10$ & MMLU & base & 0.776 & 0.754 & 1.000 \\
    MMLU $n=10$ & MMLU & explain & 0.736 & 0.787 & 0.958 \\
    MMLU $n=10$ & BBH & base & 0.511 & 0.600 & 0.542 \\
    MMLU $n=10$ & BBH & explain & 0.412 & 0.721 & 0.534 \\
    \bottomrule
  \end{tabular}
\end{table}

The cosine analysis supports a narrower conclusion than activation steering
would require. The repeated-wrong contrast produces a stable layer-23 direction
within MMLU, especially from moderate to high participant counts, and this
direction has partial but weaker alignment with BBH. Explanation prompting does
not remove the direction; rather, its directions are less aligned with the MMLU
base references on average. These results are consistent with the behavioral
pattern while remaining descriptive: they show representational similarity
between contrast directions, not causal control over the final answer.
